% ============================================================
%  Chapter 3: Literature Review
%  Included from main.tex via \input{litreview}.
%  Do NOT add \documentclass or \begin{document} here.
%  Citation style: IEEEtran.bst. All \cite keys must resolve
%  to entries in references.bib.
% ============================================================

\chapter{Literature Review}
\label{chap:litreview}

This chapter critically reviews the prior work that bears on the
dissertation's design decisions and research gap. Each section argues
against a named criterion rather than summarising papers in sequence,
and every citation is tied to a specific downstream consequence for
the design chapter. Ethical, legal, and sustainability literature is
handled in Chapter~\ref{chap:lsesp}; OCR benchmarks without a
downstream extraction component and closed-source commercial systems
without published methodology are out of scope.

% ------------------------------------------------------------
\section{The template-free criterion and what it excludes}
\label{sec:lr-templatefree}

\begin{figure}[ht]
\centering
\includegraphics[width=\linewidth]{Chapters/figures/fig-3-3-extraction-paradigm-tree.png}
\caption[Taxonomy of extraction approaches.]{Taxonomy of extraction
approaches surveyed in this chapter, organised by their answer to
the template-free criterion.}
\label{fig:lr-taxonomy}
\end{figure}

Prior work on structured extraction from semi-structured documents can
be organised around a single criterion drawn from
\S\ref{sec:bg-semistructured}: can the method process an unseen
provider, in an unseen language, with no per-document configuration or
training data? A method that fails this test cannot be deployed against
the real world of European utility bills, where the next bill in the
queue may come from a provider the system has never seen.

The rule-based and 2D-grid families surveyed in
\S\ref{sec:bg-template} fail this criterion in the same way: each
binds extraction to a learnable layout distribution, so an unseen
issuer breaks the prior. CloudScan~\cite{palm2017cloudscan} is the
clearest demonstration: on 326{,}471 invoices, the authors separate
performance on previously-seen layouts from performance on unseen
ones and report the resulting gap, establishing cross-issuer
generalisation as a distinct evaluation axis. Benchmark-trained
extractors on CORD~\cite{park2019cord}-class receipts inherit the
same problem at smaller scale. Representation learning over extraction
candidates~\cite{majumder2020form} relaxes the layout-prior
requirement but not the training-data requirement: a labelled
target-schema dataset is still needed at training time.

Specialised end-to-end document models fail the criterion for a
related reason. Donut~\cite{kim2022donut},
LayoutLMv3~\cite{huang2022layoutlmv3},
Pix2Struct~\cite{lee2023pix2struct},
Florence-2~\cite{xiao2024florence}, and
InternVL~1.5~\cite{chen2024internvl} all reach their reported accuracy
only under task-specific fine-tuning on benchmark corpora such as
DUE~\cite{borchmann2021due} and
DocVQA~\cite{mathew2021docvqachallenge}, which places them outside
the template-free zero-shot regime this project evaluates. More
fundamentally, none has a text-mode counterpart sharing its reasoning
engine, so any cross-modality comparison using them swaps the
architecture when it swaps the modality, confounding the two
variables. They serve as absolute-performance reference points but
not as the experimental vehicle for isolating input modality.

Every approach surveyed above shares one pre-condition: a layout
prior, a labelled training set in the target schema, or both. A
multilingual utility-bill corpus with providers from four countries
violates both pre-conditions by construction. Only general-purpose
zero-shot LLM and VLM prompting satisfies the template-free
criterion; which modality to feed such a model is the question the
rest of this chapter works on.

% ------------------------------------------------------------
\section{OCR-cascade pipelines and the error-propagation ceiling}
\label{sec:lr-ocrcascade}

An OCR-cascade pipeline feeds the OCR stage's output into a downstream
reasoner. Classical cascade work on the ICDAR 2019 SROIE
competition~\cite{huang2021sroie} established the evaluation framing
but tuned systems to a 1{,}000-image benchmark and did not target
cross-issuer generalisation. LLM-based cascade work relaxes the final
stage: Khanchandani et al.~\cite{khanchandani2026invoice} combine
CNN-based layout analysis, OCR, and LLM entity recognition, arguing
that an LLM reasoning stage complements the OCR layer with
contextual semantic mapping, and Abdallah et
al.~\cite{abdallah2024receiptsense} provide a
multilingual Arabic--English receipt dataset with a Tesseract OCR
baseline for LLM evaluation. Both satisfy the template-free criterion
of \S\ref{sec:lr-templatefree} because the LLM absorbs semantic
variation that would otherwise demand per-issuer configuration.

The shared ceiling of every cascade variant is OCR quality on
multilingual noisy input. Gupte et al.~\cite{gupte2021lightscameras}
synthesise controlled OCR noise and show downstream NER accuracy
drops monotonically, with a text-restoration model reducing but not
closing the gap. Hamdi et al.~\cite{hamdi2023ocrner} report that NER
F1 tracks character error rate closely on historical multilingual
text, and Nazeem et al.~\cite{nazeem2024ocrlow} document material
per-language variance across three open-source OCR engines on four
low-resource Indian languages, relevant here because a Tesseract-based
pipeline inherits whatever per-language accuracy profile Tesseract
offers.

Post-OCR correction with language models has been proposed as a
mitigation. CLOCR-C~\cite{bourne2024clocrc} reports large CER
reductions on historical OCR output; Thomas et
al.~\cite{thomas2024postocr} compare instruction-tuned Llama~2
against fine-tuned BART on the same task class. The constraint on
both is semantic: an OCR stage that has confidently misread
\texttt{177.53} as \texttt{l77.53} has produced a plausible-looking
sequence that a correction LLM has no obvious reason to change. The
cascade literature therefore establishes a ceiling on multilingual
noisy input but does not establish whether that ceiling is binding
when compared to skipping OCR entirely with the same reasoner. That
comparison requires an OCR-free counterfactual using the same
reasoning engine, which the cascade literature does not supply.

These observations feed the Failure-Attribution and
Multilingual-Coverage requirements in \S\ref{sec:lr-gap}.

% ------------------------------------------------------------
\section{Text-mode LLM extraction and the output-mechanism asymmetry}
\label{sec:lr-textmode}

A growing literature uses general-purpose LLMs (GPT-4 family, Claude,
Gemini) for structured extraction from text input, typically the
output of an OCR stage or a PDF text layer.
ChatIE~\cite{wei2023chatie} establishes the zero-shot capability with
a two-stage chat prompting scheme on relation triples, NER, and event
extraction across multiple datasets and languages, showing that
general-purpose chat LLMs perform template-free IE without
task-specific training. Biswas and Talukdar~\cite{biswas2024rotation}
extend the robustness picture under controlled in-plane rotation
across three multimodal LLMs, reporting skew-dependent hallucination
behaviour on document images.

The methodological concern for any multi-provider study is how
structured output is produced, because coverage is uneven across
providers. Geng et al.'s JSONSchemaBench~\cite{geng2025jsonbench}
evaluates six constrained-decoding frameworks (Guidance, Outlines,
Llamacpp, XGrammar, OpenAI, Gemini) and reports that efficiency,
coverage, and output quality vary meaningfully across them. OpenAI's
Structured Outputs~\cite{openai2024structured,openai-structuredoutputs-docs}
guarantee schema-valid JSON for supported models; Anthropic's Claude
does not expose an equivalent guarantee and relies on
prompt-instructed JSON with post-hoc parsing.

Structural validity is orthogonal to semantic faithfulness. OpenAI's
own documentation~\cite{openai-structuredoutputs-docs} flags that
value-level hallucinations can persist under strict schema following,
and ExtractBench~\cite{ferguson2026extractbench} pairs 35 PDFs with
JSON Schemas and explicitly separates omission errors from
hallucinations in nested-field evaluation across frontier models. A
reported accuracy difference between two providers in a text-mode
extraction study can therefore reflect a difference in reasoning, in
output-mechanism reliability, or in both, and existing studies
typically do not control for the second factor.

These observations feed the Cross-Provider-Generality and
Reproducibility requirements in \S\ref{sec:lr-gap}.

% ------------------------------------------------------------
\section{Vision-mode VLM extraction: accuracy known, trade-offs under-characterised}
\label{sec:lr-visionmode}

General-purpose VLMs accept document images directly and skip the OCR
stage. The literature on using them for structured extraction reads
along three axes: accuracy, cost, and the OCR-vs-vision tension
inside the VLM itself.

On accuracy, frontier VLMs are capable but uneven. Across the three
sources cited below, GPT-4~Turbo, GPT-4V, GPT-4o, Claude~3.5 Sonnet,
and Gemini are placed on comparable evaluation axes by Borchmann's
reproducible API evaluation of GPT-4~Turbo and GPT-4V on DocVQA and
InfoVQA~\cite{borchmann2024gpt4}, the InternVL~1.5
tables~\cite{chen2024internvl}, and the MM-Vet~v2 capability
benchmark~\cite{yu2024mmvet2}.
On multilingual and non-Latin text the picture is weaker: Fu et
al.'s bilingual benchmark~\cite{fu2024bilingual} reports that most
state-of-the-art LMMs score below 50 out of 100 on text-centric
tasks, and Shi et al.~\cite{shi2023gpt4v} find GPT-4V struggles on
multilingual OCR and complex layouts and does not outperform
specialised OCR overall.

Cost characterisations are thinner. Vasu et al.'s
FastVLM~\cite{vasu2024fastvlm} analyses the trade-off between image
resolution, vision-encoder latency, visual-token count, and LLM
inference cost, showing that document-appropriate high resolutions
sharply increase multimodal cost unless the vision encoder is
engineered for token efficiency. The practitioner question --- what
does a given level of extraction accuracy cost in tokens and dollars
at a given resolution --- is answered in parts but not end-to-end.

The analytically interesting finding comes from
DocVLM~\cite{nacson2024docvlm}, which argues that document
understanding needs fine-grained text processing \emph{and} high
resolution, that OCR text alone can underperform full-resolution
vision on dense documents, and that neither modality dominates on all
documents. On some documents the dominant information is visual
(layout, table structure, logos); on others it is textual and the OCR
layer already captures it. Whether utility bills fall consistently in
one camp or vary across issuers is exactly the empirical question
this dissertation poses.

These observations feed the Operational-Instrumentation requirement
in \S\ref{sec:lr-gap}.

% ------------------------------------------------------------
\section{Controlled comparisons of OCR-mediated vs vision-mediated extraction}
\label{sec:lr-controlled}

A small but growing literature compares OCR-cascade pipelines against
VLM vision-mode pipelines on document extraction tasks. Four controls
collectively define a satisfying answer to the practitioner question
posed here:

\begin{description}
\item[(i)] The reasoning engine is held constant across modalities.
\item[(ii)] The prompt is held constant across modalities.
\item[(iii)] The corpus covers more than one language.
\item[(iv)] Cost and latency are measured alongside accuracy.
\end{description}

\begin{figure}[ht]
\centering
\includegraphics[width=\linewidth]{Chapters/figures/fig-3-1-studies-vs-controls-matrix.png}
\caption[Controls coverage across surveyed studies.]{Coverage of the
four rubric controls across surveyed cross-modality comparisons.
No existing study meets all four; the target committed to in this
dissertation is shown in the last row.}
\label{fig:lr-controls-matrix}
\end{figure}

Each study below is assessed against these four controls. Individual
studies satisfy some and not others; no study satisfies all four.

\paragraph{The closest analogue.} Shen et al.'s 2026 EACL
industry-track paper~\cite{shen2026ocrornot} is the nearest comparator
to the present project. They benchmark out-of-the-box MLLMs on
business-document information extraction, contrasting image-only MLLM
pipelines with OCR-enhanced MLLM setups, and conclude that OCR may not
be necessary for powerful MLLMs because image-only input can reach
comparable performance. The paper satisfies controls (i) and,
arguably, (iv) within its industry framing, but does not stratify by
language or isolate the prompt as a strictly held variable. Its
central finding supplies the closest available directional prediction
for what this dissertation will measure; its controls do not answer
the multilingual strict-isolation version of the question.

\paragraph{Domain-specific comparators.} Nunes et
al.~\cite{nunes2025tables} compare Table Transformer plus OCR against
multimodal LLMs on PubTables-1M and report that CV+OCR is slightly
better on structural layout while MLLMs are stronger on cell content
--- a finding that swaps the architecture with the modality, so
control (i) fails by design. Wang and Shen~\cite{wang2025hybrid}
evaluate 25 OCR+LLM configurations on identity documents and report
$F_1 = 1.0$ on structured inputs and $F_1 = 0.997$ on challenging
images via adaptive routing, but the reasoning engine is not held
constant and the corpus is monolingual. Berghaus et
al.~\cite{berghaus2025multimodal} benchmark eight multimodal LLMs
from three families on three open invoice datasets under zero-shot
prompting, comparing direct image input against markdown-mediated
text input, but without cost/latency instrumentation or multilingual
stratification.

\paragraph{Broad-coverage evaluation suites.} Benkirane et al.'s
DISCO suite~\cite{benkirane2026disco} evaluates OCR pipelines and
VLMs separately on parsing and QA across document types including
multilingual scripts, medical forms, infographics, and multi-page
documents. OCR pipelines remain more reliable on handwriting and
long documents; VLMs perform better on multilingual text and
visually rich layouts. The complexity stratification is the
methodological strength; the limitation, for present purposes, is
that the two paradigms are evaluated as separate systems rather than
as competing input modalities for the same reasoning engine.
Controls (i) and (ii) fail.

\paragraph{Cost-focused comparisons.} olmOCR~\cite{poznanski2025olmocr}
fine-tunes a 7B-parameter VLM on a 260{,}000-page corpus and reports
comparison against GPT-4o, Gemini Flash, and Qwen-2.5-VL, with
explicit per-million-page pricing (approximately \$176 for olmOCR
against approximately \$6{,}240 for GPT-4o batch pricing on their
documents). The cost characterisation is the strongest available;
control (i) does not apply because the comparison is
cross-architecture by design.

\paragraph{Vendor benchmarks.} The Unstructured SCORE-Bench
results~\cite{unstructured2025scorebench} place certain VLM
configurations ahead of OCR configurations on the same parser, but
the reasoning engine is not held constant across rows and the
rankings are produced and maintained by the vendor.

\paragraph{Summary.} Shen et al. and DISCO ask the general question
at high abstraction and stop short of strict single-variable
isolation; the rest answer narrower questions (tables, identity
documents, invoices, cost-per-page, parsing fidelity) without holding
the reasoning engine constant. The gap is not that no one has
compared OCR and vision extraction; it is that no study
simultaneously (i) holds the reasoning engine and prompt absolutely
constant, (ii) uses a multilingual corpus from the genuinely
out-of-distribution population of European utility bills, (iii)
reports cost and latency as first-class metrics, and (iv) varies the
LLM provider factorially so the modality finding can be tested for
provider specificity.

% ------------------------------------------------------------
\section{Evaluation methodology for structured extraction}
\label{sec:lr-eval}

Even a well-controlled comparison produces uninterpretable results if
the evaluation metrics conflate distinct failure modes. This section
reviews how the extraction literature measures success against one
criterion: does the methodology distinguish semantically equivalent
predictions from genuine failures, and correct-absence from both
omission and hallucination?

MUC-7~\cite{chinchor1998muc7} supplies the foundational vocabulary:
each predicted slot is tallied as correct, incorrect, missing (MIS),
or spurious (SPU), with precision, recall, and F derived from these
counts, so under-generation and over-generation are explicitly
distinguished. Much of the field has the vocabulary available but
does not fully apply it. DocVQA's ANLS metric~\cite{mathew2021docvqa}
allows partial credit through string similarity but conflates
correct-absence with missing answers. FUNSD~\cite{jaume2019funsd}
reports entity-level precision and recall without a null taxonomy,
and DUE~\cite{borchmann2021due} aggregates per-task metrics so null
handling varies accordingly.

LMDX~\cite{perot2024lmdx} introduces a grounding-based evaluation in
which each predicted entity is tied to a source span and unverifiable
entities are discarded as hallucinations, but it does not standardise
how absent ground-truth fields are scored.
ExtractBench~\cite{ferguson2026extractbench} addresses the issue head
on: 35 PDFs paired with JSON schemas and gold labels, with omission
errors and hallucinations reported as separate categories in
nested-field evaluation across frontier models. It provides the
peer-reviewed language for the distinction this dissertation's
evaluation framework implements.

Cross-field reading from outside the extraction literature helps.
Dale et al.'s machine-translation hallucination
study~\cite{dale2023hallucination} defines hallucinations as
translations containing information unrelated to the source and
omissions as translations dropping source information. The two
concepts map cleanly onto structured extraction: a predicted field
that does not exist on the document is a hallucination; a field on
the document that the prediction omits is an omission.

Reproducibility sits alongside. Pineau et
al.~\cite{pineau2021reproducibility} describe the NeurIPS
reproducibility programme, which has since become standard for
evaluation reporting; most of the extraction literature reviewed in
this chapter does not meet the checklist in full, with config,
prompt, and per-document artefacts rarely archived.

These observations feed the Null-Aware Evaluation,
Failure-Attribution, and Reproducibility requirements in
\S\ref{sec:lr-gap}.

% ------------------------------------------------------------
\section{Research gap and bridge to specification and design}
\label{sec:lr-gap}

Across the three literatures examined --- OCR-cascade structured
extraction, text-mode LLM extraction, and cross-modality comparison
--- no published study simultaneously isolates the reasoning engine,
varies the LLM provider factorially, reports per-stage cost and
latency, and applies a null-aware evaluation taxonomy on a
multilingual out-of-distribution document class. This synthesis
insight is stated by no single paper cited above; it is the product
of reading the three literatures against one another and against the
four-point rubric of \S\ref{sec:lr-controlled}.

Based on the limitations identified in
Sections~\ref{sec:lr-ocrcascade} (OCR error propagation on
multilingual input), \ref{sec:lr-textmode} (output-mechanism
asymmetry across providers), \ref{sec:lr-controlled} (the four-point
rubric that existing comparisons fail), and \ref{sec:lr-eval}
(null-handling and failure-attribution gaps in standard evaluation),
this dissertation commits to seven experimental requirement classes
that collectively close the gap:
\textit{Modality-Isolation}, \textit{Multilingual-Coverage},
\textit{Operational-Instrumentation}, \textit{Reproducibility},
\textit{Failure-Attribution}, \textit{Null-Aware Evaluation}, and
\textit{Cross-Provider Generality}. Each class is formally stated,
given a unique identifier (ER1--ER7), and traced back to the
specific subsection of this chapter that motivates it in
Chapter~\ref{chap:spec}~\S\ref{sec:spec-experimental-requirements};
the architectural realisation of each, together with the
alternatives considered and rejected, is documented in
Chapter~\ref{chap:design}. Figure~\ref{fig:lr-gap-map}
visualises the traceability explicitly.

\begin{figure}[ht]
\centering
\includegraphics[width=\linewidth]{Chapters/figures/fig-3-2-gap-to-requirements-map.png}
\caption[Gap-to-requirement traceability map.]{Traceability from
literature gaps to the seven experimental requirement classes
formalised in Chapter~\ref{chap:spec}.}
\label{fig:lr-gap-map}
\end{figure}

Research questions RQ1--RQ3
(\S\ref{sec:intro-rqs}) operationalise the gap identified here and
are restated, paired with directional hypotheses, in
Chapter~\ref{chap:spec}~\S\ref{sec:spec-objectives}.
